\chapter{Instalación, compilación y reproducción de los experimentos}\label{anexo_instalacion}

Este anexo resume la organización del proyecto y los pasos necesarios para preparar el entorno, compilar las extensiones nativas, ejecutar las pruebas y reproducir el perfilado. El repositorio combina herramientas CMake/CUDA para la parte nativa y \texttt{uv}/Python para las implementaciones y pruebas de alto nivel.

\section{Organización del proyecto}

El flujo de trabajo se divide en los siguientes componentes:

\begin{itemize}
    \item \textbf{CMake}: configura y compila el código C++/CUDA, las extensiones nativas y los tests de bajo nivel.
    \item \textbf{uv}: crea el entorno Python reproducible e instala las dependencias declaradas por el proyecto.
    \item \textbf{pytest}: ejecuta el test funcional, los benchmarks temporales y los tests preparados para \textit{profiling}.
    \item \textbf{CTest}: ejecuta desde CMake los tests correspondientes a la parte C++/CUDA.
    \item \textbf{\texttt{profile.sh}}: automatiza las sesiones de NVIDIA Nsight Compute para CUDA, Triton, SDPA y SDPA-cuDNN.
    \item \textbf{\path{environment_report.sh}}: registra las características de hardware y las versiones de software relevantes para reproducir las mediciones.
\end{itemize}

Los artefactos de compilación se almacenan bajo el directorio \texttt{build/}. Los informes generados por Nsight Compute se almacenan en \texttt{profile/}, de forma que los resultados de perfilado quedan separados del código fuente y de los binarios compilados.

\section{Requisitos}

La configuración utilizada en la memoria emplea CUDA Toolkit 13.0, CMake, un compilador C++ compatible con C++20, Python gestionado mediante \texttt{uv}, PyTorch, Triton y NVIDIA Nsight Compute. La implementación CUDA final está especializada para la arquitectura Ampere y el preset suministrado configura como objetivo \texttt{sm\_80}, correspondiente a la A100 utilizada en las mediciones.

El repositorio incluye un preset CMake denominado \texttt{angel-wsl}. En dicho preset se especifican, entre otros parámetros, la arquitectura CUDA objetivo y las rutas de los compiladores. Estas rutas dependen de la máquina concreta, por lo que el usuario puede modificar el preset para adaptarlo a su instalación. De la misma forma, si se utiliza otra GPU Ampere deberá revisarse la arquitectura indicada en \texttt{CUDA\_ARCHITECTURES}.

\section{Preparación del entorno Python}

El entorno y las dependencias se sincronizan mediante:

\begin{lstlisting}[language=bash]
uv sync
\end{lstlisting}

Una vez completado este paso, las herramientas Python del proyecto pueden ejecutarse mediante \texttt{uv run}, evitando depender de paquetes instalados globalmente en el sistema.

\section{Configuración y compilación con CMake}

La primera configuración del árbol de compilación puede realizarse con el preset incluido:

\begin{lstlisting}[language=bash]
cmake --preset angel-wsl
\end{lstlisting}

La compilación se realiza mediante:

\begin{lstlisting}[language=bash]
cmake --build --preset angel-wsl
\end{lstlisting}

Los ficheros generados permanecen dentro de \texttt{build/}. Si se modifica la ruta de CUDA, el compilador C++ o la arquitectura objetivo, debe actualizarse el preset antes de regenerar el proyecto.

\section{Ejecución de tests}

La batería Python completa puede prepararse y ejecutarse con:

\begin{lstlisting}[language=bash]
uv sync && uv run pytest
\end{lstlisting}

Los tests de rendimiento utilizan la misma semilla pseudoaleatoria en todas las implementaciones. Concretamente, antes de generar \(Q\), \(K\) y \(V\) se inicializan los generadores con semilla 0 mediante \path{torch.manual_seed(0)} y \path{torch.cuda.manual_seed_all(0)}. De esta manera, CUDA, Triton, SDPA y cuDNN se evalúan con entradas reproducibles generadas bajo el mismo criterio, evitando que diferencias en los datos aleatorios introduzcan una fuente adicional de variación entre benchmarks.

Los tests C++/CUDA integrados con CMake pueden ejecutarse mediante CTest desde el árbol de compilación:

\begin{lstlisting}[language=bash]
ctest --test-dir build --output-on-failure
\end{lstlisting}

Si el preset utiliza un subdirectorio específico dentro de \texttt{build/}, el argumento de \texttt{--test-dir} debe apuntar al directorio generado por dicho preset.

\section{Perfilado con Nsight Compute}

El script \texttt{profile.sh} automatiza la captura de las regiones NVTX correspondientes a las cuatro rutas evaluadas. Para habilitar su ejecución y lanzar el perfilado se utiliza:

\begin{lstlisting}[language=bash]
chmod u+x ./profile.sh && ./profile.sh
\end{lstlisting}

El script genera informes \texttt{.ncu-rep} en el directorio \texttt{profile/}. Estos informes pueden abrirse posteriormente con \texttt{ncu-ui} para analizar ocupación, registros, memoria compartida, tráfico de memoria, conflictos de banco, utilización de Tensor Cores y causas de espera de los \textit{warps}.

Las duraciones obtenidas bajo Nsight Compute no se utilizan como benchmark temporal principal, ya que la herramienta puede instrumentar y repetir internamente el kernel. Las medidas de tiempo presentadas en la memoria proceden de los tests específicos basados en eventos CUDA.

\section{Registro del entorno de ejecución}

Para facilitar la reproducibilidad se incluye el script \path{environment_report.sh}. Su ejecución se realiza mediante:

\begin{lstlisting}[language=bash,breaklines=true,columns=fullflexible]
chmod u+x ./environment_report.sh && ./environment_report.sh
\end{lstlisting}

El informe recoge información como sistema operativo, CPU visible, memoria del sistema, modelo de GPU, número de SM, capacidad de cómputo, frecuencias observadas, límite de potencia, driver NVIDIA, CUDA Toolkit, compiladores, CMake, Python, PyTorch, cuDNN, Triton y versión de Nsight Compute. Este procedimiento fue el utilizado para documentar el entorno de prueba descrito en la memoria.

\section{Flujo recomendado para reproducir el trabajo}

Un flujo completo de reproducción puede resumirse en los siguientes pasos:

\begin{enumerate}
    \item Adaptar, si fuese necesario, el preset \texttt{angel-wsl} a las rutas de compilador y arquitectura CUDA de la máquina.
    \item Ejecutar \texttt{uv sync} para crear el entorno Python.
    \item Configurar y compilar mediante CMake.
    \item Ejecutar \texttt{uv run pytest} y CTest para comprobar la correctitud antes de medir rendimiento.
    \item Ejecutar \path{environment_report.sh} para registrar las condiciones del experimento.
    \item Ejecutar los benchmarks temporales de \texttt{pytest} bajo las mismas condiciones y semilla.
    \item Ejecutar \texttt{profile.sh} para generar los informes de Nsight Compute en \texttt{profile/}.
\end{enumerate}

Esta separación entre compilación, validación, benchmark y perfilado evita mezclar los costes de instrumentación con las medidas de tiempo y facilita repetir el análisis en otra máquina modificando únicamente los parámetros dependientes del entorno.
